%%%%%%%%%%%%%%%%%%%%%%%%%%%%%%%%%%%%%%%%%%%%%%%%%%%%%%%%%%%%%%%%%%%%%%
% Lớp 12 | Ứng dụng đạo hàm để khảo sát hàm số | Khảo sát sự biến thiên và vẽ đồ thị hàm số
% Dạng 2D1?5-6: Phương trình tiếp tuyến của đồ thị hàm số
% Số câu: 7  (sinh tự động bởi tach_cau.py; từ nay thầy sửa trực tiếp file này)
%%%%%%%%%%%%%%%%%%%%%%%%%%%%%%%%%%%%%%%%%%%%%%%%%%%%%%%%%%%%%%%%%%%%%%

%%%%%%%%%%%%%%%% TRẮC NGHIỆM 4 PHƯƠNG ÁN %%%%%%%%%%%%%%%%

% nguồn: Phan1_Lop_12_Dai_so_Chuong_1
\begin{ex}%[2D1N5-6]%[ma:2806087a]
	Đạo hàm của hàm số $y= f(x)$ tại điểm $x_0$ là hệ số góc của tiếp tuyến với đồ thị $(C)$ của hàm số tại điểm $M_0(x_0; f(x_0))$. Khi đó phương trình tiếp tuyến của $(C)$ tại điểm $M_0$  là
	\choice 
		{$y+y_0=f'(x_0)\cdot  (x-x_0)$}
		{$y-y_0=f'(x_0)\cdot  (x+x_0)$}
		{\True $y-y_0=f'(x_0)\cdot  (x-x_0)$}
		{$y-y_0=f'(x_0)\cdot  (x_0-x)$}
	\loigiai{ 
	 }
\end{ex}

% nguồn: Phan1_Lop_12_Dai_so_Chuong_1
\begin{ex}%[2D1H5-6]%[ma:2b6ed5e8]
	Cho hàm số $y= x^3-2x+ 1.$ Viết phương trình  tiếp tuyến của đồ thị hàm số tại điểm $M(0;1)$.
	 \choice
		{$ y=2x  +3$}
		{\True $ y=-2x + 1$} 
		{$ y=4x + 1$} 
		{$ y=-4x + 1$} 
	\loigiai{ 
		Đạo hàm của hàm số đã cho là $y'= 3x^2-2$ $\Rightarrow$ $y'(0)=-2$.\\ 
		Phương trình tiếp tuyến của đồ thị hàm số tại điểm $M(0;1)$  là\\
		$y-1=  -2(x-0)$ hay $ y=-2x + 1$.
	} 
\end{ex}


%%%%%%%%%%%%%%%% ĐÚNG/SAI %%%%%%%%%%%%%%%%

% nguồn: Phan1_Lop_12_Dai_so_Chuong_1
\begin{ex}%[2D1H5-6]%[ma:3f4aba5b]
	 Cho hàm số $f(x)=x^3-x+1$.
	\choiceTF 
		 {\True $f'(x)=3x^2-1$}
		 {\True $f'(1)=2$}
		 {$f(1)=0$}
		{Phương trình tiếp tuyến của đồ thị  hàm số đã cho tại điểm có hoành độ $x_{0}=1$ là $y=2x-2$}
	 \loigiai{
		\begin{itemchoice}
			\itemch Đúng. 
			\itemch  Đúng.
			 \itemch Sai. Vì $f(1) =1^3-1+1 =1 \ne 0$.
			\itemch Sai. Vì phương trình tiếp tuyến  của hàm số $y=f(x)$ tại $x_0$ có dạng \[y=f'(x_0)(x-x_0) + f(x_0).\]
			Áp dụng vào hàm số $y=x^3-x+1$ ta có $f(1)=1$,  $f'(1)=2$, thay vào công~thức trên ta có $y=2(x-1)+1 = 2x-1$.
		\end{itemchoice}
	 }
\end{ex}

% nguồn: Phan1_Lop_12_Dai_so_Chuong_1
\begin{ex}%[2D1H5-6]%[ma:c5f6c616]
	Cho  hàm số $y=\dfrac{-mx^2+(4m-2)x+1-4m}{x-1}$ có đồ thị là $(C)$ với $m$ là tham số.
	 \choiceTF
		{\True Khi $m=1$ đồ thị hàm số  có 2 điểm cực trị}
		{\True Khi $m=1$ đồ thị hàm số không cắt  trục $Ox$}
		{\True Khi $m<-1$ thì hàm số có  điểm cực đại và điểm cực tiểu trong miền $x>0$}
		{\True Khi $m=1$ thì có  2 phương trình tiếp tuyến của $(C)$ song song với đường thẳng $x-y=0$}
	 \loigiai{
		\begin{itemchoice}
			\itemch 
			 Khi $m=1$ ta có $y=\dfrac{-x^2+2x-3}{x-1}=-x+1-\dfrac{2}{x-1}$.\\
			Tập  xác định: $\mathbb{D}=\mathbb{R}\setminus\{1\}$;\\
			 Ta có $y'=\dfrac{-x^2+2x+1}{(x-1)^2}=0\Leftrightarrow-x^2+2x+1=0\Leftrightarrow\hoac{&x=1-\sqrt{2}\Rightarrow y=2\sqrt{2}\\&x=1+\sqrt{2}\Rightarrow y=-2\sqrt{2}}$.
			Ta lại có $\lim\limits_{x\to 1^{\pm}}y=\pm\infty$ nên $x=1$  là tiệm cận đứng; $\lim\limits_{x\to\pm\infty}y=-x+1$ nên $y=-x+1$ là tiệm cận xiên.
			Bảng biến thiên: 
			\begin{center}\begin{tikzpicture}[>=stealth]
		\tkzTabInit[nocadre=false,lgt=1,espcl=2,deltacl=0.5]{$x$/.7,$y'$/.7,$y$/2} 
			{$-\infty$,$1-\sqrt{2}$,$1$,$1+\sqrt{2}$,$+\infty$} 
		 \tkzTabLine{,-,$0$,+,d,+,$0$,-,}
		\tkzTabVar{+/$+\infty$,-/$2\sqrt{2}$,+D-/$+\infty$/$-\infty$,+/$-2\sqrt{2}$,-/$-\infty$} 
	\end{tikzpicture}
			\end{center}
			Vậy khi $m=1$ đồ thị hàm số có 2 điểm cực  trị.
			\itemch Khi $m=1$ ta có  $y=\dfrac{-x^2+2x-3}{x-1}=-x+1-\dfrac{2}{x-1}$.\\
			Tập xác định:  $\mathbb{D}=\mathbb{R}\setminus\{1\}$;\\
			$y=0\Rightarrow-x^2+2x-3=0$ (vô nghiệm) nên đồ  thị hàm số không cắt trục $Ox$.
			 \itemch $y=\dfrac{-mx^2+(4m-2)x+1-4m}{x-1}\Rightarrow y'=\dfrac{-mx^2+2mx-4m+2-1+4m}{(x-1)^2}$.\\
			Dấu $y'$ là  dấu của tam thức $g(x) = -mx^2 + 2mx + 1$.\\
			$g(x)$  có $\Delta'=m^2+m$; $g(1)=-m+2m+1=m+1$
			Để hàm số có cực đại và cực tiểu  thì
			$\heva{& \Delta'>0 \\&m + 1 \ne 0}\Leftrightarrow\hoac{&m<-1\\&m>0}$ (1). 
			Lúc này hàm số đạt cực đại và cực tiểu  tại $x=x_1$, $x=x_2$ và
			$\heva{&x_1 + x_2 =  2 \\&x_1\cdot x_2 = -\dfrac{1}{m}.}$
			Theo yêu  cầu bài toán $\heva{&x_1>0\\&x_2>0}\Leftrightarrow
			\heva{&x_1 +  x_2 > 0 \\&x_1 \cdot x_2 > 0}\Leftrightarrow m < 0$ (2)
			Từ (1) và (2) ta có $m < -1$. 
			\itemch Khi $m=1$ thì $y'_x = \dfrac{-x^2 + 2x +  1}{(x-1)^2}$. Đường thẳng $x - y = 0$ có hệ số góc $k = 1$.
			Để tiếp tuyến của $(C)$ song song với  đường thẳng $y = x$ cần là $y'(x) = 1$ \\
			 $\Leftrightarrow \heva{&x \ne 1 \\ &\dfrac{-x^2 + 2x + 1}{(x - 1)^2} = 1 } \Leftrightarrow \heva{&x \ne 1 \\ &-x^2 + 2x + 1 = (x - 1)^2 } \Leftrightarrow \heva{&x \ne 1 \\ &-x^2 + 2x + 1 = x^2 - 2x + 1 }$\\
			$ \Leftrightarrow \heva{&x \ne 1 \\ &-2x^2 + 4x = 0 } \Leftrightarrow \hoac{&x = 0 \Rightarrow y  = 3. \\&x = 2 \Rightarrow y = -3.
			 }$
			Có hai tiếp tuyến thỏa yêu  cầu bài toán:
			$(T_1): y = 1(x-0) + 3 \Leftrightarrow y  = x + 3$;\quad $(T_2): y = 1(x-2) - 3 \Leftrightarrow y = x - 5$.
		\end{itemchoice}
	 }
\end{ex}


%%%%%%%%%%%%%%%% TRẢ LỜI NGẮN %%%%%%%%%%%%%%%%

% nguồn: Phan2_Lop_12_Dai_so_Chuong_1
\begin{bt}%[2D1V5-6]%[ma:450bca2f]
	Cho hàm số $y=\dfrac{1}{3}x^{3}-2x^{2}+3x+5$. Biết phương trình tiếp tuyến  của đồ thị hàm số với hệ số góc nhỏ nhất có dạng là $y=ax+b \left( a,b\in \mathbb{R} \right)$. Tính giá trị của biểu thức $T=a+3b$.
	\shortans{22} 
	 \loigiai{
		 Ta có: $y'=x^{2}-4x+3$.\\
		 Khi đó, $\min y'=y'\left( 2 \right)=-1\Rightarrow x^{2}-4x+3=-1\Rightarrow x=2\Rightarrow y=\dfrac{17}{3}$.\\
		Phương trình tiếp  tuyến là $y=-1\left( x-2 \right)+\dfrac{17}{3}=-x+\dfrac{23}{3}$. \\
		 Vậy $T=-1+3\cdot \dfrac{23}{3}=22$.
	} 
\end{bt}

% nguồn: Phan2_Lop_12_Dai_so_Chuong_1
\begin{bt}%[2D1V5-6]%[ma:c243bd5d]
	\immini{Một đường ray tàu lượn trong khu vui chơi giải trí có hình dáng được  mô phỏng theo đồ thị của hàm số $y=x^3-3x+2$, ký hiệu là $(C)$. Để đảm bảo an toàn và tính thẩm mỹ, người ta chọn hai điểm $A(a;b)$ và $B(c;d)$ trên đường ray sao cho tiếp tuyến tại hai điểm này có cùng độ dốc (cùng hệ số góc). Đồng thời, đoạn đường nối hai trụ đỡ tại các điểm $A$ và $B$ phải vuông góc với một
	đường dây điện có phương trình $x+y-5=0$.  Tìm $b+d$.}{
	\begin{tikzpicture}[line cap=butt,line join=miter,>=stealth]
		\draw[->] (-3,0)--(3,0) 
			 node[shift={(-100:7pt)},font=\normalsize]{$x$};
		\draw[->]  (0,-1)--(0,5)
			node[shift={(180:7pt)},font=\normalsize]{$y$}; 
		\draw (5pt,0) |- (0,5pt) 
			(0,0)  node[shift={(225:9pt)},font=\normalsize]{$O$};
		\fill (0,0)  circle (1pt);
		\tikzset{declare  function={f (\x)=(\x)^3-3*(\x)+2;}}
		\begin{scope}
			\draw[blue,samples=100] plot[domain=-2.1:1.8] (\x,{f  (\x)});
		\end{scope}
	\end{tikzpicture}
	 }
	 \shortans{4}
	 \loigiai{Ta có $y=x^{3}-3x+2\Rightarrow y'=3x^{2}-3$.\\
		Tiếp tuyến với $(C)$ tại $A$, $B$ có cùng hệ số góc và chỉ khi 
		 \[f'(x_A)=f'(x_B)
		\Leftrightarrow  x_A^2=x_B^2
		\Leftrightarrow \hoac{&x_A=x_B\,(\text{loại}) \\&x_A+x_B=0.}\] 
		Suy ra $A$, $B$ đối xứng với nhau qua điểm $I(0;2)$ là tâm đối xứng của  $(C)$.\\
		$AB\perp  d\colon x+y-5=0\Rightarrow AB\colon x-y+m=0$.\\
		$AB$ qua  $I$ nên ta có $m=2\Rightarrow AB\colon x-y+2=0$.\\
		Khi đó hoành độ  $A$, $B$ thỏa mãn phương trình
		$x^{3}-3x+2=x+2\Leftrightarrow \hoac{&x=0\,(\text{loại})  \\
		 &x=\pm 2.}$\\
		$\Rightarrow  A(2;4)$, $B(-2;0)$.\\
		Do đó, $b=4$  và $d=0$.\\
		Vậy $b+d=4+0=4$. 
	 }
\end{bt}


%%%%%%%%%%%%%%%% TỰ LUẬN %%%%%%%%%%%%%%%%

% nguồn: Phan2_Lop_12_Dai_so_Chuong_1
\begin{bt}%[2D1H5-6]%[ma:9e200e79]
	Cho đường cong $(C):y={{x}^{3}}-4$. Viết phương trình tiếp tuyến của đường cong $\left( C \right)$ đi qua  điểm $A\left( 2;4 \right)$.
	 \loigiai{
		Ta có $y'=3x^2$. Gọi $M\left( x_0;y_0 \right)$ là tọa độ  tiếp điểm.\\
		Phương trình tiếp tuyến của $\left( C \right)$ tại điểm $M$ có dạng $y=3x_0^2\cdot\left( x-x_0  \right)+x_0^3-4$.\\
		Vì tiếp tuyến đi qua điểm $A\left( 2;4 \right)$ nên ta  có
		\[4=3x_0^2\cdot\left( 2-x_0 \right)+x_0^3-4\Leftrightarrow -2x_0^3+6x_0^2-8=0\Leftrightarrow\hoac{&x_0=-1\\&x_0=2.}  \]
		\begin{itemize}
			\item Với $x_0=-1$  suy ra phương trình tiếp tuyến là $y=3x-2.$
			\item Với $x_0=2$ suy ra  phương trình tiếp tuyến là $y=12x-20.$
		\end{itemize}
	} 
\end{bt}
